\scanpage{45}
\diagram{\begin{tikzcd}[ampersand replacement=\&,column sep=small]
0\arrow[r] \& \underline{\Phi}_X\arrow[r] \& \underline{\mathrm{QF}}_X\arrow[r,very thick] \& \underline{\mathrm{Div}}_X\arrow[r,dash] \& 0
\end{tikzcd}}
\diagram{\begin{tikzcd}[ampersand replacement=\&,column sep=small]
0\arrow[r] \& \underline R_X^*/\underline U_X\arrow[r] \& \underline{\mathrm{QF}}_X\arrow[r] \& \underline{\Phi}_X/\operatorname{Im}\underline O_X^*\arrow[r] \& 0
\end{tikzcd}}

Qu'on a avantage à envisager en diagramme \unsure{commun|comme|commutatif}.

\diagram{\footnotesize\begin{tikzcd}[ampersand replacement=\&,column sep=small,row sep=1.4em]
\& \& 0 \& \& \\
0\arrow[r,phantom,"{\simeq}" description] \& \underline{\Phi}_X/\underline{\operatorname{Im}}\underline O_X^*\arrow[r,phantom,"{\simeq}" description] \& \underline{\Phi}_X/\underline{\operatorname{Im}}\underline O_X^*\arrow[u]\arrow[r,phantom,"{\simeq}" description] \& 0 \& \\
0\arrow[r,dash] \& \underline{\Phi}_X\arrow[u,dashed]\arrow[r] \& \underline{\mathrm{QF}}_X\arrow[u]\arrow[r] \& \underline{\mathrm{Div}}_X\arrow[u,dash,dashed,"{\unsure{\sim|?|\wr}}" sloped] \& \\
0\arrow[r] \& \underline O_X^*/\underline U_X\arrow[u]\arrow[r] \& \underline R_X^*/\underline U_X\arrow[u]\arrow[r,dashed] \& \underline{\mathrm{Div}}_X\arrow[u,dashed,"{\sim}" sloped]\arrow[r,dashed] \& 0\\
\& 0\arrow[u] \& 0\arrow[u] \& 0\arrow[u,dash,dashed,"{\sim}" sloped] \&
\end{tikzcd}}

\unsure{,|.|{}} Passant aux sections globales, on trouve $\bigl\{$ pour $X$ géom intègre [ce qui implique que $\underline R_X^*/\underline U_X$ est constant, donc flasque $\bigr\}$, on obtient le diagramme commutatif à lignes et colonnes exactes
% archive sheet number: 44
